\documentclass[Total.tex]{subfiles}
\begin{document}
	The first part is the basic theory of Banach space. This is an extra contents for the functional analysis. 
	
	\chapter{Banach Spaces}
		\section{Definitions, and Basic Properties}
			\begin{Def}
				A normed space $(X,\Vert-\Vert)$ is called a \textbf{Banach space} if it is compelete as a metric space. 
			\end{Def}
			\subsection{Banach Space from Normed Space: Compeletion}
				\begin{proposition}
					内容...
				\end{proposition}
				\begin{proof}
					内容...
				\end{proof}
				
	\chapter{Hilbert Spaces}
		Recall:
		
		\begin{lemma}
			(Jordan-von Neumann)TFAE:
			\begin{itemize}
				\item 
				\item 
			\end{itemize}
		\end{lemma}
		\begin{proof}
				内容...
			\end{proof}
		
		
		\section{Pre-Hilbert Space}
		\section{Hilbert Space}
			\begin{Def}
				A \textbf{Hilbert space} is a $F$-linear space which is complete in the metric induced by the norm.
				
				i.e. it is a Banach space. 
			\end{Def}

			According to the equivalence : Norm $\Leftrightarrow$ Inner Product.	

			\begin{proposition}
				Let $H$ be not complete with respect to its induced norm, then its Banach space compeletion is a Hilbert space.
			\end{proposition}
			\begin{proof}
				
			\end{proof}
			
			\subsection{Duality and Representation}
			\subsubsection{Riesz's Representation Theorem}
			
			\subsubsection{Self-Duality}
			
			\subsubsection{Reflxivity}
			
			\subsubsection{Riesz-Frechet}
			
			\subsubsection{Lax-Milgram Theorem}
			
		\chapter{Operators on Hilbert Spaces}
			
	\newpage The second part is the Banach space theory.
		
	\newpage The third part is the Interpolation Theory.
	
		\chapter{Classical Interpolation Theory}
		
		\chapter{Absolute Interpolation Theory}
\end{document}